% !TeX program = pdflatex
% Komentarz do umowy spółki Basilisk Systems P.S.A. (dokument informacyjny, odrębny od umowy).
% Numery paragrafów pobierane z psa.aux (xr-hyper); bibliografia wspólna: basilisk_bib.tex.
\documentclass[11pt,a4paper]{report}
\usepackage[T1]{fontenc}
\usepackage[utf8]{inputenc}
\usepackage[english,polish]{babel}
\usepackage{xr-hyper}
\usepackage{basilisk_i18n}
\usepackage{lmodern}
\usepackage{microtype}
\usepackage{setspace}
\usepackage{parskip}
\usepackage{needspace}
\externaldocument{psa}[psa.pdf]

\BusinessPlanCompany{Basilisk Systems}
\BusinessPlanWatermarkText{Basilisk Systems}
\BusinessPlanWatermarkColor{black!6}
\BusinessPlanWatermarkFontSize{38}{46}
\BusinessPlanWatermarkAngle{38}
\BusinessPlanAuthorsText{Basilisk Systems --- materiały pomocnicze}
\BusinessPlanConfidentialText{Dokument informacyjny --- nie stanowi części umowy}
\BusinessPlanFooterText{Basilisk Systems --- komentarz do umowy spółki}
\BusinessPlanVersionText{Wersja 1.0 --- projekt}

\setstretch{1.08}
\setlength{\parindent}{0pt}
\setlength{\parskip}{0.5em}
\setlength{\emergencystretch}{4em}
\renewcommand{\arraystretch}{1.18}

\makeatletter
\newcommand{\BibItem}[2]{\item[\textbf{[#1]}]\def\@currentlabel{#1}\def\@currentlabelname{#1}\phantomsection\label{bib:#2}}
\makeatother
\newcommand{\Zr}[1]{[\ref{bib:#1}]}
\newcommand{\U}[1]{\textsuperscript{#1}}
% wiersz nagłówka paragrafu i wiersz ustępu
\newcommand{\Par}[2]{\multicolumn{3}{@{}l}{\rule{0pt}{1.4em}\textbf{§~\ref{art:#1}. #2}} \\ \midrule}
\newcommand{\Ust}[3]{#1 & #2 & #3 \\}

\hypersetup{
  pdftitle={Basilisk Systems P.S.A. --- komentarz do umowy spółki},
  pdfauthor={Basilisk Systems}
}

\begin{document}

\BusinessPlanTitlePage
  {Basilisk Systems}
  {Wersja 1.0 --- projekt}
  {Komentarz do umowy spółki}
  {Podstawy prawne i bibliografia do umowy Basilisk Systems P.S.A.}
  {Do umowy spółki w wersji 3.3-S --- per paragraf i ustęp}
  {1.0}

\tableofcontents
\newpage

\chapter*{Wprowadzenie}
\addcontentsline{toc}{chapter}{Wprowadzenie}

Dokument jest \textbf{informacyjny}: nie stanowi części umowy spółki, nie jest objęty aktem notarialnym ani podpisem i nie zmienia treści umowy. Służy założycielom, dyrektorom, przyszłym inwestorom i ich doradcom do szybkiego odnalezienia podstawy prawnej każdego postanowienia. W razie rozbieżności rozstrzyga treść umowy spółki \Zr{us}.

Dokument składa się z dwóch załączników informacyjnych:
\begin{itemize}
  \item \textbf{Załącznik informacyjny A --- Bibliografia}: akty prawne krajowe i unijne, orzecznictwo, literatura, materiały urzędowe i dokumenty Spółki; etykiety w nawiasach kwadratowych są klikalne;
  \item \textbf{Załącznik informacyjny B --- Podstawy prawne per paragraf i ustęp}: dla każdego ustępu umowy przepisy, na których się opiera, orzecznictwo i krótki komentarz. Numery paragrafów są klikalne i otwierają umowę (\texttt{psa.pdf} w tym samym katalogu).
\end{itemize}

Stan prawny: listopad 2026 --- według tekstu jednolitego KSH (Dz.U. 2024 poz. 18 ze zm., z uwzględnieniem nowelizacji Dz.U. 2026 poz. 176, która od 18 lutego 2027 r. zmienia przepisy o rejestrze akcjonariuszy) i pozostałych aktów z bibliografii. Numery przepisów należy zweryfikować w ISAP przed użyciem w obrocie; tezy orzeczeń podano w skrócie. Skrót „PSA” oznacza prostą spółkę akcyjną, „WZ” --- Walne Zgromadzenie.

% ================================================================
\chapter*{Załącznik informacyjny A}
\addcontentsline{toc}{chapter}{Załącznik informacyjny A --- Bibliografia}
% Wspólna bibliografia Basilisk Systems --- % Wspólna bibliografia Basilisk Systems --- % Wspólna bibliografia Basilisk Systems --- \input{basilisk_bib} w podręczniku decyzji i komentarzu do umowy spółki.
% Odesłanie w tekście: \Zr{klucz} -> [ETYKIETA] z linkiem do pozycji.
\makeatletter
\providecommand{\BibItem}[2]{\item[\textbf{[#1]}]\def\@currentlabel{#1}\def\@currentlabelname{#1}\phantomsection\label{bib:#2}}
\makeatother
\providecommand{\Zr}[1]{[\ref{bib:#1}]}
\newenvironment{BibList}{\begin{list}{}{\setlength{\labelwidth}{2.6cm}\setlength{\leftmargin}{2.9cm}\setlength{\labelsep}{0.3cm}\setlength{\itemsep}{0.15em}\setlength{\parsep}{0pt}\small}}{\end{list}}

\chapter*{Bibliografia}
\addcontentsline{toc}{chapter}{Bibliografia}

Akty prawne podano według miejsca pierwotnej publikacji oraz ostatniego tekstu jednolitego (stan: listopad 2026), z dopiskiem „ze zm.”; przed użyciem należy sprawdzić, czy nie ogłoszono nowszego tekstu w Internetowym Systemie Aktów Prawnych (isap.sejm.gov.pl) albo w Dzienniku Ustaw (eli.gov.pl). Tezy orzeczeń podano w skrócie; przed powołaniem się na nie w sporze należy sprawdzić pełny tekst.

\section*{A. Akty prawne krajowe}
\begin{BibList}
\BibItem{KSH}{ksh} Ustawa z dnia 15 września 2000 r. --- Kodeks spółek handlowych (Dz.U. 2000 nr 94 poz. 1037 ze zm.; t.j. Dz.U. 2024 poz. 18 ze zm., w tym ustawą z dnia 23 stycznia 2026 r., Dz.U. 2026 poz. 176 --- zmiany rejestru akcjonariuszy od 18 lutego 2027 r.); przepisy o prostej spółce akcyjnej: dział IA, art.~300$^{1}$--300$^{134}$.
\BibItem{Nowela PSA}{psa2019} Ustawa z dnia 19 lipca 2019 r. o zmianie ustawy --- Kodeks spółek handlowych oraz niektórych innych ustaw (Dz.U. 2019 poz. 1655) --- wprowadzenie prostej spółki akcyjnej; ustawa w pierwotnym brzmieniu wchodziła w życie 1 marca 2020 r., a przepisy o PSA stosuje się od 1 lipca 2021 r. (po przesunięciu terminu).
\BibItem{KC}{kc} Ustawa z dnia 23 kwietnia 1964 r. --- Kodeks cywilny (Dz.U. 1964 nr 16 poz. 93; t.j. Dz.U. 2026 poz. 795 ze zm.).
\BibItem{KRO}{kro} Ustawa z dnia 25 lutego 1964 r. --- Kodeks rodzinny i opiekuńczy (Dz.U. 1964 nr 9 poz. 59; t.j. Dz.U. 2026 poz. 236 ze zm.).
\BibItem{KPC}{kpc} Ustawa z dnia 17 listopada 1964 r. --- Kodeks postępowania cywilnego (Dz.U. 1964 nr 43 poz. 296; t.j. Dz.U. 2026 poz. 468 ze zm.).
\BibItem{KP}{kp} Ustawa z dnia 26 czerwca 1974 r. --- Kodeks pracy (Dz.U. 1974 nr 24 poz. 141; t.j. Dz.U. 2025 poz. 277 ze zm.).
\BibItem{PrAut}{praut} Ustawa z dnia 4 lutego 1994 r. o prawie autorskim i prawach pokrewnych (Dz.U. 1994 nr 24 poz. 83; t.j. Dz.U. 2025 poz. 24 ze zm.).
\BibItem{PWP}{pwp} Ustawa z dnia 30 czerwca 2000 r. --- Prawo własności przemysłowej (Dz.U. 2001 nr 49 poz. 508; t.j. Dz.U. 2023 poz. 1170 ze zm.).
\BibItem{UZNK}{uznk} Ustawa z dnia 16 kwietnia 1993 r. o zwalczaniu nieuczciwej konkurencji (Dz.U. 1993 nr 47 poz. 211; t.j. Dz.U. 2026 poz. 85 ze zm.).
\BibItem{PIT}{pit} Ustawa z dnia 26 lipca 1991 r. o podatku dochodowym od osób fizycznych (Dz.U. 1991 nr 80 poz. 350; t.j. Dz.U. 2026 poz. 592 ze zm.).
\BibItem{CIT}{cit} Ustawa z dnia 15 lutego 1992 r. o podatku dochodowym od osób prawnych (Dz.U. 1992 nr 21 poz. 86 ze zm.).
\BibItem{UoR}{uor} Ustawa z dnia 29 września 1994 r. o rachunkowości (Dz.U. 1994 nr 121 poz. 591; t.j. Dz.U. 2026 poz. 522 ze zm.).
\BibItem{UKRS}{ukrs} Ustawa z dnia 20 sierpnia 1997 r. o Krajowym Rejestrze Sądowym (Dz.U. 1997 nr 121 poz. 769; t.j. Dz.U. 2025 poz. 869 ze zm.).
\BibItem{UONI}{uoni} Ustawa z dnia 5 sierpnia 2010 r. o ochronie informacji niejawnych (Dz.U. 2010 nr 182 poz. 1228; t.j. Dz.U. 2025 poz. 1209 ze zm.).
\BibItem{UOTS}{uots} Ustawa z dnia 29 listopada 2000 r. o obrocie z zagranicą towarami, technologiami i usługami o znaczeniu strategicznym dla bezpieczeństwa państwa, a także dla utrzymania międzynarodowego pokoju i bezpieczeństwa (Dz.U. 2000 nr 119 poz. 1250; t.j. Dz.U. 2023 poz. 1582 ze zm., w tym Dz.U. 2026 poz. 471 --- definicje i zezwolenia, m.in. pomoc techniczna).
\BibItem{UWT}{uwt} Ustawa z dnia 13 czerwca 2019 r. o wykonywaniu działalności gospodarczej w zakresie wytwarzania i obrotu materiałami wybuchowymi, bronią, amunicją oraz wyrobami i technologią o przeznaczeniu wojskowym lub policyjnym (Dz.U. 2019 poz. 1214; t.j. Dz.U. 2023 poz. 1743 ze zm.); koncesja na obrót technologią o przeznaczeniu wojskowym (art.~7), wymogi wobec wspólników z co najmniej 20\% akcji (art.~10 i 17)..
\BibItem{UKI}{uki} Ustawa z dnia 24 lipca 2015 r. o kontroli niektórych inwestycji (Dz.U. 2015 poz. 1272; t.j. Dz.U. 2026 poz. 47 ze zm.).
\BibItem{USank}{usank} Ustawa z dnia 13 kwietnia 2022 r. o szczególnych rozwiązaniach w zakresie przeciwdziałania wspieraniu agresji na Ukrainę oraz służących ochronie bezpieczeństwa narodowego (Dz.U. 2022 poz. 835; t.j. Dz.U. 2025 poz. 514 ze zm.).
\BibItem{USyg}{usyg} Ustawa z dnia 14 czerwca 2024 r. o ochronie sygnalistów (Dz.U. 2024 poz. 928).
\BibItem{UODO}{uodo} Ustawa z dnia 10 maja 2018 r. o ochronie danych osobowych (Dz.U. 2018 poz. 1000 ze zm.).
\BibItem{PrLot}{prlot} Ustawa z dnia 3 lipca 2002 r. --- Prawo lotnicze (Dz.U. 2002 nr 130 poz. 1112; t.j. Dz.U. 2025 poz. 1431 ze zm.).
\BibItem{PKD 2025}{pkd} Rozporządzenie Rady Ministrów z dnia 18 grudnia 2024 r. w sprawie Polskiej Klasyfikacji Działalności (PKD) (Dz.U. 2024 poz. 1936).
\BibItem{WT}{wt} Rozporządzenie Rady Ministrów z dnia 17 września 2019 r. w sprawie klasyfikacji rodzajów materiałów wybuchowych, broni, amunicji oraz wyrobów i technologii o przeznaczeniu wojskowym lub policyjnym, na których wytwarzanie lub obrót jest wymagane uzyskanie koncesji (Dz.U. 2019 poz. 1888) --- obejmuje m.in. bezzałogowe statki powietrzne i ich systemy (WT~V pkt~3).
\end{BibList}

\section*{B. Akty prawne Unii Europejskiej}
\begin{BibList}
\BibItem{RODO}{rodo} Rozporządzenie Parlamentu Europejskiego i Rady (UE) 2016/679 z dnia 27 kwietnia 2016 r. w sprawie ochrony osób fizycznych w związku z przetwarzaniem danych osobowych i w sprawie swobodnego przepływu takich danych oraz uchylenia dyrektywy 95/46/WE (ogólne rozporządzenie o ochronie danych) (Dz.Urz. UE L 119 z 4.5.2016, s.~1).
\BibItem{Dual-use}{dualuse} Rozporządzenie Parlamentu Europejskiego i Rady (UE) 2021/821 z dnia 20 maja 2021 r. ustanawiające unijny system kontroli wywozu, pośrednictwa, pomocy technicznej, tranzytu i transferu produktów podwójnego zastosowania (wersja przekształcona) (Dz.Urz. UE L 206 z 11.6.2021, s.~1, ze zm.).
\BibItem{EDF}{edf} Rozporządzenie Parlamentu Europejskiego i Rady (UE) 2021/697 z dnia 29 kwietnia 2021 r. ustanawiające Europejski Fundusz Obronny i uchylające rozporządzenie (UE) 2018/1092 (Dz.Urz. UE L 170 z 12.5.2021, s.~149).
\BibItem{UAS 945}{uas945} Rozporządzenie delegowane Komisji (UE) 2019/945 z dnia 12 marca 2019 r. w sprawie bezzałogowych systemów powietrznych oraz operatorów bezzałogowych systemów powietrznych z państw trzecich (Dz.Urz. UE L 152 z 11.6.2019, s.~1).
\BibItem{UAS 947}{uas947} Rozporządzenie wykonawcze Komisji (UE) 2019/947 z dnia 24 maja 2019 r. w sprawie przepisów i procedur dotyczących eksploatacji bezzałogowych statków powietrznych (Dz.Urz. UE L 152 z 11.6.2019, s.~45).
\BibItem{FDI}{fdi} Rozporządzenie Parlamentu Europejskiego i Rady (UE) 2019/452 z dnia 19 marca 2019 r. ustanawiające ramy monitorowania bezpośrednich inwestycji zagranicznych w Unii --- uchylone z dniem 17 stycznia 2028 r. przez \Zr{fdi2026}.
\BibItem{FDI 2026}{fdi2026} Rozporządzenie Parlamentu Europejskiego i Rady (UE) 2026/1386 z dnia 17 czerwca 2026 r. w sprawie monitorowania inwestycji zagranicznych w Unii i uchylające rozporządzenie (UE) 2019/452 (Dz.Urz. UE L 2026/1386 z 26.6.2026); stosuje się od 17 stycznia 2028 r. (niektóre przepisy od 16 lipca 2026 r.); obejmuje m.in. spółki rozwijające produkty podwójnego zastosowania, wojskowe i drony.
\BibItem{TTBER}{ttber} Rozporządzenie Komisji (UE) 2026/877 z dnia 16 kwietnia 2026 r. w sprawie stosowania art.~101 ust.~3 TFUE do kategorii porozumień o transferze technologii (Dz.Urz. UE L 2026/877 z 21.4.2026); obowiązuje od 1 maja 2026 r. do 30 kwietnia 2038 r.; zastąpiło rozporządzenie (UE) nr 316/2014 --- licencje na technologię, know-how i oprogramowanie.
\BibItem{AI Act}{aiact} Rozporządzenie Parlamentu Europejskiego i Rady (UE) 2024/1689 z dnia 13 czerwca 2024 r. w sprawie ustanowienia zharmonizowanych przepisów dotyczących sztucznej inteligencji (akt w sprawie sztucznej inteligencji) (Dz.Urz. UE L 2024/1689 z 12.7.2024); systemy wyłącznie wojskowe są wyłączone (art.~2 ust.~3), systemy podwójnego zastosowania --- objęte.
\BibItem{Sankcje}{sankcje} Rozporządzenie Rady (UE) nr 833/2014 z dnia 31 lipca 2014 r. dotyczące środków ograniczających w związku z działaniami Rosji destabilizującymi sytuację na Ukrainie (Dz.Urz. UE L 229 z 31.7.2014, s.~1) oraz rozporządzenie Rady (UE) nr 269/2014 z dnia 17 marca 2014 r. (Dz.Urz. UE L 78 z 17.3.2014, s.~6) (środki ograniczające).
\BibItem{TSD}{tsd} Dyrektywa Parlamentu Europejskiego i Rady (UE) 2016/943 z dnia 8 czerwca 2016 r. w sprawie ochrony niejawnego know-how i niejawnych informacji handlowych (tajemnic przedsiębiorstwa) przed ich bezprawnym pozyskiwaniem, wykorzystywaniem i ujawnianiem (Dz.Urz. UE L 157 z 15.6.2016, s.~1).
\BibItem{WBD}{wbd} Dyrektywa Parlamentu Europejskiego i Rady (UE) 2019/1937 z dnia 23 października 2019 r. w sprawie ochrony osób zgłaszających naruszenia prawa Unii.
\end{BibList}

\section*{C. Orzecznictwo}
\begin{BibList}
\BibItem{SN 61/03}{sn6103} Uchwała składu siedmiu sędziów SN z dnia 6 listopada 2003 r., III CZP 61/03, OSNC 2004, nr 5, poz. 69 (zasada prawna) --- kara umowna należy się także wtedy, gdy wierzyciel nie poniósł szkody; dłużnik może żądać jej miarkowania (art.~484 §~2 KC).
\BibItem{SN 304/00}{sn30400} Wyrok SN z dnia 3 października 2000 r., I CKN 304/00, OSNC 2001, nr 4, poz. 59 --- wykorzystanie przez pracownika informacji, co do których przedsiębiorca nie podjął niezbędnych działań w celu zachowania ich poufności, należy traktować jako wykorzystanie powszechnej wiedzy, do której przedsiębiorca nie ma ustawowych uprawnień.
\BibItem{SN 120/07}{sn12007} Wyrok SN z dnia 8 stycznia 2008 r., II PK 120/07 --- odpowiedzialność pracownika za szkodę regulują wyczerpująco przepisy Kodeksu pracy; kara umowna (art.~483 KC) zastrzeżona w stosunku pracy nie może być dochodzona; dopuszczalna jest natomiast w umowie o zakazie konkurencji po ustaniu stosunku pracy (SN powołuje wyrok z 10 października 2003 r., I PK 528/02, OSNP 2004, nr 19, poz. 336).
\BibItem{SN 42/18}{sn4218} Uchwała SN z dnia 7 września 2018 r., III CZP 42/18 (Lex nr 2541937) --- członek zarządu nie może jako pełnomocnik reprezentować spółki z o.o. w umowie z innym członkiem zarządu (art.~210 §~1 KSH); podobnie uchwała SN z 30 stycznia 2019 r., III CZP 71/18. Skutkiem naruszenia jest bezwzględna nieważność czynności (art.~58 §~1 KC) --- tak wyrok SA w Warszawie z 9 kwietnia 2019 r., V ACa 485/18, który obie uchwały przywołuje. Przydatna przy wykładni art.~300$^{79}$ KSH.
\end{BibList}

\section*{D. Literatura}
\begin{BibList}
\BibItem{Malinowski}{malinowski} P. Malinowski (red.), \emph{Prosta spółka akcyjna. Komentarz}, Warszawa 2021.
\BibItem{Opalski}{opalski} A. Opalski, \emph{Prosta spółka akcyjna --- nowy typ spółki handlowej}, cz.~I: „Przegląd Prawa Handlowego” 2019, nr 11, s.~5--15; cz.~II: tamże, nr 12, s.~4--16.
\BibItem{ECFR}{ecfr} \emph{The Simple Joint Stock Company: Emergence of a New Close Company in Poland}, „European Company and Financial Law Review” 2024, t.~21, DOI 10.1515/ecfr-2024-0021.
\BibItem{Jara}{jara} Z. Jara (red.), \emph{Kodeks spółek handlowych. Komentarz}, C.H. Beck (wydanie aktualne).
\BibItem{Gniewek}{gniewek} E. Gniewek, P. Machnikowski (red.), \emph{Kodeks cywilny. Komentarz}, C.H. Beck (wydanie aktualne).
\BibItem{Barta}{barta} J. Barta, R. Markiewicz (red.), \emph{Ustawa o prawie autorskim i prawach pokrewnych. Komentarz}, Wolters Kluwer (wydanie aktualne).
\BibItem{Kostański}{kostanski} P. Kostański (red.), \emph{Prawo własności przemysłowej. Komentarz}, C.H. Beck (wydanie aktualne).
\BibItem{Szwaja}{szwaja} J. Szwaja (red.), \emph{Ustawa o zwalczaniu nieuczciwej konkurencji. Komentarz}, C.H. Beck (wydanie aktualne).
\end{BibList}

\section*{E. Materiały urzędowe}
\begin{BibList}
\BibItem{GUS PKD}{guspkd} Główny Urząd Statystyczny, \emph{Polska Klasyfikacja Działalności PKD 2025} (klasyfikacje.stat.gov.pl).
\BibItem{GUS CPI}{guscpi} Komunikaty Prezesa GUS w sprawie średniorocznego wskaźnika cen towarów i usług konsumpcyjnych ogółem, ogłaszane w Dzienniku Urzędowym RP „Monitor Polski” (podstawa waloryzacji kar i Cennika).
\end{BibList}

\section*{F. Dokumenty Spółki}
\begin{BibList}
\BibItem{US}{us} Umowa spółki Basilisk Systems P.S.A., wersja 3.3-S (\texttt{psa.pdf}; wersja angielska \texttt{psa\_en.pdf}).
\BibItem{UA}{ua} Umowa akcjonariuszy, wersja 1.7 (\texttt{basilisk\_sha.pdf}).
\BibItem{IP}{ip} Umowy przeniesienia praw autorskich (Umowa nr~1 i nr~2), wersja 1.6 (\texttt{basilisk\_ip.pdf}).
\BibItem{Cennik}{cennik} Cennik wynagrodzeń za rozwiązania projektowe (\texttt{basilisk\_cennik.pdf}).
\BibItem{NDA}{nda} Wzór umowy o zachowaniu poufności (\texttt{basilisk\_nda.pdf}).
\BibItem{Polityka}{polityka} Polityka bezpieczeństwa informacji (\texttt{basilisk\_polityka.pdf}).
\BibItem{Podręcznik}{podrecznik} Podręcznik podejmowania decyzji (\texttt{basilisk\_decyzje.pdf}).
\end{BibList}
 w podręczniku decyzji i komentarzu do umowy spółki.
% Odesłanie w tekście: \Zr{klucz} -> [ETYKIETA] z linkiem do pozycji.
\makeatletter
\providecommand{\BibItem}[2]{\item[\textbf{[#1]}]\def\@currentlabel{#1}\def\@currentlabelname{#1}\phantomsection\label{bib:#2}}
\makeatother
\providecommand{\Zr}[1]{[\ref{bib:#1}]}
\newenvironment{BibList}{\begin{list}{}{\setlength{\labelwidth}{2.6cm}\setlength{\leftmargin}{2.9cm}\setlength{\labelsep}{0.3cm}\setlength{\itemsep}{0.15em}\setlength{\parsep}{0pt}\small}}{\end{list}}

\chapter*{Bibliografia}
\addcontentsline{toc}{chapter}{Bibliografia}

Akty prawne podano według miejsca pierwotnej publikacji oraz ostatniego tekstu jednolitego (stan: listopad 2026), z dopiskiem „ze zm.”; przed użyciem należy sprawdzić, czy nie ogłoszono nowszego tekstu w Internetowym Systemie Aktów Prawnych (isap.sejm.gov.pl) albo w Dzienniku Ustaw (eli.gov.pl). Tezy orzeczeń podano w skrócie; przed powołaniem się na nie w sporze należy sprawdzić pełny tekst.

\section*{A. Akty prawne krajowe}
\begin{BibList}
\BibItem{KSH}{ksh} Ustawa z dnia 15 września 2000 r. --- Kodeks spółek handlowych (Dz.U. 2000 nr 94 poz. 1037 ze zm.; t.j. Dz.U. 2024 poz. 18 ze zm., w tym ustawą z dnia 23 stycznia 2026 r., Dz.U. 2026 poz. 176 --- zmiany rejestru akcjonariuszy od 18 lutego 2027 r.); przepisy o prostej spółce akcyjnej: dział IA, art.~300$^{1}$--300$^{134}$.
\BibItem{Nowela PSA}{psa2019} Ustawa z dnia 19 lipca 2019 r. o zmianie ustawy --- Kodeks spółek handlowych oraz niektórych innych ustaw (Dz.U. 2019 poz. 1655) --- wprowadzenie prostej spółki akcyjnej; ustawa w pierwotnym brzmieniu wchodziła w życie 1 marca 2020 r., a przepisy o PSA stosuje się od 1 lipca 2021 r. (po przesunięciu terminu).
\BibItem{KC}{kc} Ustawa z dnia 23 kwietnia 1964 r. --- Kodeks cywilny (Dz.U. 1964 nr 16 poz. 93; t.j. Dz.U. 2026 poz. 795 ze zm.).
\BibItem{KRO}{kro} Ustawa z dnia 25 lutego 1964 r. --- Kodeks rodzinny i opiekuńczy (Dz.U. 1964 nr 9 poz. 59; t.j. Dz.U. 2026 poz. 236 ze zm.).
\BibItem{KPC}{kpc} Ustawa z dnia 17 listopada 1964 r. --- Kodeks postępowania cywilnego (Dz.U. 1964 nr 43 poz. 296; t.j. Dz.U. 2026 poz. 468 ze zm.).
\BibItem{KP}{kp} Ustawa z dnia 26 czerwca 1974 r. --- Kodeks pracy (Dz.U. 1974 nr 24 poz. 141; t.j. Dz.U. 2025 poz. 277 ze zm.).
\BibItem{PrAut}{praut} Ustawa z dnia 4 lutego 1994 r. o prawie autorskim i prawach pokrewnych (Dz.U. 1994 nr 24 poz. 83; t.j. Dz.U. 2025 poz. 24 ze zm.).
\BibItem{PWP}{pwp} Ustawa z dnia 30 czerwca 2000 r. --- Prawo własności przemysłowej (Dz.U. 2001 nr 49 poz. 508; t.j. Dz.U. 2023 poz. 1170 ze zm.).
\BibItem{UZNK}{uznk} Ustawa z dnia 16 kwietnia 1993 r. o zwalczaniu nieuczciwej konkurencji (Dz.U. 1993 nr 47 poz. 211; t.j. Dz.U. 2026 poz. 85 ze zm.).
\BibItem{PIT}{pit} Ustawa z dnia 26 lipca 1991 r. o podatku dochodowym od osób fizycznych (Dz.U. 1991 nr 80 poz. 350; t.j. Dz.U. 2026 poz. 592 ze zm.).
\BibItem{CIT}{cit} Ustawa z dnia 15 lutego 1992 r. o podatku dochodowym od osób prawnych (Dz.U. 1992 nr 21 poz. 86 ze zm.).
\BibItem{UoR}{uor} Ustawa z dnia 29 września 1994 r. o rachunkowości (Dz.U. 1994 nr 121 poz. 591; t.j. Dz.U. 2026 poz. 522 ze zm.).
\BibItem{UKRS}{ukrs} Ustawa z dnia 20 sierpnia 1997 r. o Krajowym Rejestrze Sądowym (Dz.U. 1997 nr 121 poz. 769; t.j. Dz.U. 2025 poz. 869 ze zm.).
\BibItem{UONI}{uoni} Ustawa z dnia 5 sierpnia 2010 r. o ochronie informacji niejawnych (Dz.U. 2010 nr 182 poz. 1228; t.j. Dz.U. 2025 poz. 1209 ze zm.).
\BibItem{UOTS}{uots} Ustawa z dnia 29 listopada 2000 r. o obrocie z zagranicą towarami, technologiami i usługami o znaczeniu strategicznym dla bezpieczeństwa państwa, a także dla utrzymania międzynarodowego pokoju i bezpieczeństwa (Dz.U. 2000 nr 119 poz. 1250; t.j. Dz.U. 2023 poz. 1582 ze zm., w tym Dz.U. 2026 poz. 471 --- definicje i zezwolenia, m.in. pomoc techniczna).
\BibItem{UWT}{uwt} Ustawa z dnia 13 czerwca 2019 r. o wykonywaniu działalności gospodarczej w zakresie wytwarzania i obrotu materiałami wybuchowymi, bronią, amunicją oraz wyrobami i technologią o przeznaczeniu wojskowym lub policyjnym (Dz.U. 2019 poz. 1214; t.j. Dz.U. 2023 poz. 1743 ze zm.); koncesja na obrót technologią o przeznaczeniu wojskowym (art.~7), wymogi wobec wspólników z co najmniej 20\% akcji (art.~10 i 17)..
\BibItem{UKI}{uki} Ustawa z dnia 24 lipca 2015 r. o kontroli niektórych inwestycji (Dz.U. 2015 poz. 1272; t.j. Dz.U. 2026 poz. 47 ze zm.).
\BibItem{USank}{usank} Ustawa z dnia 13 kwietnia 2022 r. o szczególnych rozwiązaniach w zakresie przeciwdziałania wspieraniu agresji na Ukrainę oraz służących ochronie bezpieczeństwa narodowego (Dz.U. 2022 poz. 835; t.j. Dz.U. 2025 poz. 514 ze zm.).
\BibItem{USyg}{usyg} Ustawa z dnia 14 czerwca 2024 r. o ochronie sygnalistów (Dz.U. 2024 poz. 928).
\BibItem{UODO}{uodo} Ustawa z dnia 10 maja 2018 r. o ochronie danych osobowych (Dz.U. 2018 poz. 1000 ze zm.).
\BibItem{PrLot}{prlot} Ustawa z dnia 3 lipca 2002 r. --- Prawo lotnicze (Dz.U. 2002 nr 130 poz. 1112; t.j. Dz.U. 2025 poz. 1431 ze zm.).
\BibItem{PKD 2025}{pkd} Rozporządzenie Rady Ministrów z dnia 18 grudnia 2024 r. w sprawie Polskiej Klasyfikacji Działalności (PKD) (Dz.U. 2024 poz. 1936).
\BibItem{WT}{wt} Rozporządzenie Rady Ministrów z dnia 17 września 2019 r. w sprawie klasyfikacji rodzajów materiałów wybuchowych, broni, amunicji oraz wyrobów i technologii o przeznaczeniu wojskowym lub policyjnym, na których wytwarzanie lub obrót jest wymagane uzyskanie koncesji (Dz.U. 2019 poz. 1888) --- obejmuje m.in. bezzałogowe statki powietrzne i ich systemy (WT~V pkt~3).
\end{BibList}

\section*{B. Akty prawne Unii Europejskiej}
\begin{BibList}
\BibItem{RODO}{rodo} Rozporządzenie Parlamentu Europejskiego i Rady (UE) 2016/679 z dnia 27 kwietnia 2016 r. w sprawie ochrony osób fizycznych w związku z przetwarzaniem danych osobowych i w sprawie swobodnego przepływu takich danych oraz uchylenia dyrektywy 95/46/WE (ogólne rozporządzenie o ochronie danych) (Dz.Urz. UE L 119 z 4.5.2016, s.~1).
\BibItem{Dual-use}{dualuse} Rozporządzenie Parlamentu Europejskiego i Rady (UE) 2021/821 z dnia 20 maja 2021 r. ustanawiające unijny system kontroli wywozu, pośrednictwa, pomocy technicznej, tranzytu i transferu produktów podwójnego zastosowania (wersja przekształcona) (Dz.Urz. UE L 206 z 11.6.2021, s.~1, ze zm.).
\BibItem{EDF}{edf} Rozporządzenie Parlamentu Europejskiego i Rady (UE) 2021/697 z dnia 29 kwietnia 2021 r. ustanawiające Europejski Fundusz Obronny i uchylające rozporządzenie (UE) 2018/1092 (Dz.Urz. UE L 170 z 12.5.2021, s.~149).
\BibItem{UAS 945}{uas945} Rozporządzenie delegowane Komisji (UE) 2019/945 z dnia 12 marca 2019 r. w sprawie bezzałogowych systemów powietrznych oraz operatorów bezzałogowych systemów powietrznych z państw trzecich (Dz.Urz. UE L 152 z 11.6.2019, s.~1).
\BibItem{UAS 947}{uas947} Rozporządzenie wykonawcze Komisji (UE) 2019/947 z dnia 24 maja 2019 r. w sprawie przepisów i procedur dotyczących eksploatacji bezzałogowych statków powietrznych (Dz.Urz. UE L 152 z 11.6.2019, s.~45).
\BibItem{FDI}{fdi} Rozporządzenie Parlamentu Europejskiego i Rady (UE) 2019/452 z dnia 19 marca 2019 r. ustanawiające ramy monitorowania bezpośrednich inwestycji zagranicznych w Unii --- uchylone z dniem 17 stycznia 2028 r. przez \Zr{fdi2026}.
\BibItem{FDI 2026}{fdi2026} Rozporządzenie Parlamentu Europejskiego i Rady (UE) 2026/1386 z dnia 17 czerwca 2026 r. w sprawie monitorowania inwestycji zagranicznych w Unii i uchylające rozporządzenie (UE) 2019/452 (Dz.Urz. UE L 2026/1386 z 26.6.2026); stosuje się od 17 stycznia 2028 r. (niektóre przepisy od 16 lipca 2026 r.); obejmuje m.in. spółki rozwijające produkty podwójnego zastosowania, wojskowe i drony.
\BibItem{TTBER}{ttber} Rozporządzenie Komisji (UE) 2026/877 z dnia 16 kwietnia 2026 r. w sprawie stosowania art.~101 ust.~3 TFUE do kategorii porozumień o transferze technologii (Dz.Urz. UE L 2026/877 z 21.4.2026); obowiązuje od 1 maja 2026 r. do 30 kwietnia 2038 r.; zastąpiło rozporządzenie (UE) nr 316/2014 --- licencje na technologię, know-how i oprogramowanie.
\BibItem{AI Act}{aiact} Rozporządzenie Parlamentu Europejskiego i Rady (UE) 2024/1689 z dnia 13 czerwca 2024 r. w sprawie ustanowienia zharmonizowanych przepisów dotyczących sztucznej inteligencji (akt w sprawie sztucznej inteligencji) (Dz.Urz. UE L 2024/1689 z 12.7.2024); systemy wyłącznie wojskowe są wyłączone (art.~2 ust.~3), systemy podwójnego zastosowania --- objęte.
\BibItem{Sankcje}{sankcje} Rozporządzenie Rady (UE) nr 833/2014 z dnia 31 lipca 2014 r. dotyczące środków ograniczających w związku z działaniami Rosji destabilizującymi sytuację na Ukrainie (Dz.Urz. UE L 229 z 31.7.2014, s.~1) oraz rozporządzenie Rady (UE) nr 269/2014 z dnia 17 marca 2014 r. (Dz.Urz. UE L 78 z 17.3.2014, s.~6) (środki ograniczające).
\BibItem{TSD}{tsd} Dyrektywa Parlamentu Europejskiego i Rady (UE) 2016/943 z dnia 8 czerwca 2016 r. w sprawie ochrony niejawnego know-how i niejawnych informacji handlowych (tajemnic przedsiębiorstwa) przed ich bezprawnym pozyskiwaniem, wykorzystywaniem i ujawnianiem (Dz.Urz. UE L 157 z 15.6.2016, s.~1).
\BibItem{WBD}{wbd} Dyrektywa Parlamentu Europejskiego i Rady (UE) 2019/1937 z dnia 23 października 2019 r. w sprawie ochrony osób zgłaszających naruszenia prawa Unii.
\end{BibList}

\section*{C. Orzecznictwo}
\begin{BibList}
\BibItem{SN 61/03}{sn6103} Uchwała składu siedmiu sędziów SN z dnia 6 listopada 2003 r., III CZP 61/03, OSNC 2004, nr 5, poz. 69 (zasada prawna) --- kara umowna należy się także wtedy, gdy wierzyciel nie poniósł szkody; dłużnik może żądać jej miarkowania (art.~484 §~2 KC).
\BibItem{SN 304/00}{sn30400} Wyrok SN z dnia 3 października 2000 r., I CKN 304/00, OSNC 2001, nr 4, poz. 59 --- wykorzystanie przez pracownika informacji, co do których przedsiębiorca nie podjął niezbędnych działań w celu zachowania ich poufności, należy traktować jako wykorzystanie powszechnej wiedzy, do której przedsiębiorca nie ma ustawowych uprawnień.
\BibItem{SN 120/07}{sn12007} Wyrok SN z dnia 8 stycznia 2008 r., II PK 120/07 --- odpowiedzialność pracownika za szkodę regulują wyczerpująco przepisy Kodeksu pracy; kara umowna (art.~483 KC) zastrzeżona w stosunku pracy nie może być dochodzona; dopuszczalna jest natomiast w umowie o zakazie konkurencji po ustaniu stosunku pracy (SN powołuje wyrok z 10 października 2003 r., I PK 528/02, OSNP 2004, nr 19, poz. 336).
\BibItem{SN 42/18}{sn4218} Uchwała SN z dnia 7 września 2018 r., III CZP 42/18 (Lex nr 2541937) --- członek zarządu nie może jako pełnomocnik reprezentować spółki z o.o. w umowie z innym członkiem zarządu (art.~210 §~1 KSH); podobnie uchwała SN z 30 stycznia 2019 r., III CZP 71/18. Skutkiem naruszenia jest bezwzględna nieważność czynności (art.~58 §~1 KC) --- tak wyrok SA w Warszawie z 9 kwietnia 2019 r., V ACa 485/18, który obie uchwały przywołuje. Przydatna przy wykładni art.~300$^{79}$ KSH.
\end{BibList}

\section*{D. Literatura}
\begin{BibList}
\BibItem{Malinowski}{malinowski} P. Malinowski (red.), \emph{Prosta spółka akcyjna. Komentarz}, Warszawa 2021.
\BibItem{Opalski}{opalski} A. Opalski, \emph{Prosta spółka akcyjna --- nowy typ spółki handlowej}, cz.~I: „Przegląd Prawa Handlowego” 2019, nr 11, s.~5--15; cz.~II: tamże, nr 12, s.~4--16.
\BibItem{ECFR}{ecfr} \emph{The Simple Joint Stock Company: Emergence of a New Close Company in Poland}, „European Company and Financial Law Review” 2024, t.~21, DOI 10.1515/ecfr-2024-0021.
\BibItem{Jara}{jara} Z. Jara (red.), \emph{Kodeks spółek handlowych. Komentarz}, C.H. Beck (wydanie aktualne).
\BibItem{Gniewek}{gniewek} E. Gniewek, P. Machnikowski (red.), \emph{Kodeks cywilny. Komentarz}, C.H. Beck (wydanie aktualne).
\BibItem{Barta}{barta} J. Barta, R. Markiewicz (red.), \emph{Ustawa o prawie autorskim i prawach pokrewnych. Komentarz}, Wolters Kluwer (wydanie aktualne).
\BibItem{Kostański}{kostanski} P. Kostański (red.), \emph{Prawo własności przemysłowej. Komentarz}, C.H. Beck (wydanie aktualne).
\BibItem{Szwaja}{szwaja} J. Szwaja (red.), \emph{Ustawa o zwalczaniu nieuczciwej konkurencji. Komentarz}, C.H. Beck (wydanie aktualne).
\end{BibList}

\section*{E. Materiały urzędowe}
\begin{BibList}
\BibItem{GUS PKD}{guspkd} Główny Urząd Statystyczny, \emph{Polska Klasyfikacja Działalności PKD 2025} (klasyfikacje.stat.gov.pl).
\BibItem{GUS CPI}{guscpi} Komunikaty Prezesa GUS w sprawie średniorocznego wskaźnika cen towarów i usług konsumpcyjnych ogółem, ogłaszane w Dzienniku Urzędowym RP „Monitor Polski” (podstawa waloryzacji kar i Cennika).
\end{BibList}

\section*{F. Dokumenty Spółki}
\begin{BibList}
\BibItem{US}{us} Umowa spółki Basilisk Systems P.S.A., wersja 3.3-S (\texttt{psa.pdf}; wersja angielska \texttt{psa\_en.pdf}).
\BibItem{UA}{ua} Umowa akcjonariuszy, wersja 1.7 (\texttt{basilisk\_sha.pdf}).
\BibItem{IP}{ip} Umowy przeniesienia praw autorskich (Umowa nr~1 i nr~2), wersja 1.6 (\texttt{basilisk\_ip.pdf}).
\BibItem{Cennik}{cennik} Cennik wynagrodzeń za rozwiązania projektowe (\texttt{basilisk\_cennik.pdf}).
\BibItem{NDA}{nda} Wzór umowy o zachowaniu poufności (\texttt{basilisk\_nda.pdf}).
\BibItem{Polityka}{polityka} Polityka bezpieczeństwa informacji (\texttt{basilisk\_polityka.pdf}).
\BibItem{Podręcznik}{podrecznik} Podręcznik podejmowania decyzji (\texttt{basilisk\_decyzje.pdf}).
\end{BibList}
 w podręczniku decyzji i komentarzu do umowy spółki.
% Odesłanie w tekście: \Zr{klucz} -> [ETYKIETA] z linkiem do pozycji.
\makeatletter
\providecommand{\BibItem}[2]{\item[\textbf{[#1]}]\def\@currentlabel{#1}\def\@currentlabelname{#1}\phantomsection\label{bib:#2}}
\makeatother
\providecommand{\Zr}[1]{[\ref{bib:#1}]}
\newenvironment{BibList}{\begin{list}{}{\setlength{\labelwidth}{2.6cm}\setlength{\leftmargin}{2.9cm}\setlength{\labelsep}{0.3cm}\setlength{\itemsep}{0.15em}\setlength{\parsep}{0pt}\small}}{\end{list}}

\chapter*{Bibliografia}
\addcontentsline{toc}{chapter}{Bibliografia}

Akty prawne podano według miejsca pierwotnej publikacji oraz ostatniego tekstu jednolitego (stan: listopad 2026), z dopiskiem „ze zm.”; przed użyciem należy sprawdzić, czy nie ogłoszono nowszego tekstu w Internetowym Systemie Aktów Prawnych (isap.sejm.gov.pl) albo w Dzienniku Ustaw (eli.gov.pl). Tezy orzeczeń podano w skrócie; przed powołaniem się na nie w sporze należy sprawdzić pełny tekst.

\section*{A. Akty prawne krajowe}
\begin{BibList}
\BibItem{KSH}{ksh} Ustawa z dnia 15 września 2000 r. --- Kodeks spółek handlowych (Dz.U. 2000 nr 94 poz. 1037 ze zm.; t.j. Dz.U. 2024 poz. 18 ze zm., w tym ustawą z dnia 23 stycznia 2026 r., Dz.U. 2026 poz. 176 --- zmiany rejestru akcjonariuszy od 18 lutego 2027 r.); przepisy o prostej spółce akcyjnej: dział IA, art.~300$^{1}$--300$^{134}$.
\BibItem{Nowela PSA}{psa2019} Ustawa z dnia 19 lipca 2019 r. o zmianie ustawy --- Kodeks spółek handlowych oraz niektórych innych ustaw (Dz.U. 2019 poz. 1655) --- wprowadzenie prostej spółki akcyjnej; ustawa w pierwotnym brzmieniu wchodziła w życie 1 marca 2020 r., a przepisy o PSA stosuje się od 1 lipca 2021 r. (po przesunięciu terminu).
\BibItem{KC}{kc} Ustawa z dnia 23 kwietnia 1964 r. --- Kodeks cywilny (Dz.U. 1964 nr 16 poz. 93; t.j. Dz.U. 2026 poz. 795 ze zm.).
\BibItem{KRO}{kro} Ustawa z dnia 25 lutego 1964 r. --- Kodeks rodzinny i opiekuńczy (Dz.U. 1964 nr 9 poz. 59; t.j. Dz.U. 2026 poz. 236 ze zm.).
\BibItem{KPC}{kpc} Ustawa z dnia 17 listopada 1964 r. --- Kodeks postępowania cywilnego (Dz.U. 1964 nr 43 poz. 296; t.j. Dz.U. 2026 poz. 468 ze zm.).
\BibItem{KP}{kp} Ustawa z dnia 26 czerwca 1974 r. --- Kodeks pracy (Dz.U. 1974 nr 24 poz. 141; t.j. Dz.U. 2025 poz. 277 ze zm.).
\BibItem{PrAut}{praut} Ustawa z dnia 4 lutego 1994 r. o prawie autorskim i prawach pokrewnych (Dz.U. 1994 nr 24 poz. 83; t.j. Dz.U. 2025 poz. 24 ze zm.).
\BibItem{PWP}{pwp} Ustawa z dnia 30 czerwca 2000 r. --- Prawo własności przemysłowej (Dz.U. 2001 nr 49 poz. 508; t.j. Dz.U. 2023 poz. 1170 ze zm.).
\BibItem{UZNK}{uznk} Ustawa z dnia 16 kwietnia 1993 r. o zwalczaniu nieuczciwej konkurencji (Dz.U. 1993 nr 47 poz. 211; t.j. Dz.U. 2026 poz. 85 ze zm.).
\BibItem{PIT}{pit} Ustawa z dnia 26 lipca 1991 r. o podatku dochodowym od osób fizycznych (Dz.U. 1991 nr 80 poz. 350; t.j. Dz.U. 2026 poz. 592 ze zm.).
\BibItem{CIT}{cit} Ustawa z dnia 15 lutego 1992 r. o podatku dochodowym od osób prawnych (Dz.U. 1992 nr 21 poz. 86 ze zm.).
\BibItem{UoR}{uor} Ustawa z dnia 29 września 1994 r. o rachunkowości (Dz.U. 1994 nr 121 poz. 591; t.j. Dz.U. 2026 poz. 522 ze zm.).
\BibItem{UKRS}{ukrs} Ustawa z dnia 20 sierpnia 1997 r. o Krajowym Rejestrze Sądowym (Dz.U. 1997 nr 121 poz. 769; t.j. Dz.U. 2025 poz. 869 ze zm.).
\BibItem{UONI}{uoni} Ustawa z dnia 5 sierpnia 2010 r. o ochronie informacji niejawnych (Dz.U. 2010 nr 182 poz. 1228; t.j. Dz.U. 2025 poz. 1209 ze zm.).
\BibItem{UOTS}{uots} Ustawa z dnia 29 listopada 2000 r. o obrocie z zagranicą towarami, technologiami i usługami o znaczeniu strategicznym dla bezpieczeństwa państwa, a także dla utrzymania międzynarodowego pokoju i bezpieczeństwa (Dz.U. 2000 nr 119 poz. 1250; t.j. Dz.U. 2023 poz. 1582 ze zm., w tym Dz.U. 2026 poz. 471 --- definicje i zezwolenia, m.in. pomoc techniczna).
\BibItem{UWT}{uwt} Ustawa z dnia 13 czerwca 2019 r. o wykonywaniu działalności gospodarczej w zakresie wytwarzania i obrotu materiałami wybuchowymi, bronią, amunicją oraz wyrobami i technologią o przeznaczeniu wojskowym lub policyjnym (Dz.U. 2019 poz. 1214; t.j. Dz.U. 2023 poz. 1743 ze zm.); koncesja na obrót technologią o przeznaczeniu wojskowym (art.~7), wymogi wobec wspólników z co najmniej 20\% akcji (art.~10 i 17)..
\BibItem{UKI}{uki} Ustawa z dnia 24 lipca 2015 r. o kontroli niektórych inwestycji (Dz.U. 2015 poz. 1272; t.j. Dz.U. 2026 poz. 47 ze zm.).
\BibItem{USank}{usank} Ustawa z dnia 13 kwietnia 2022 r. o szczególnych rozwiązaniach w zakresie przeciwdziałania wspieraniu agresji na Ukrainę oraz służących ochronie bezpieczeństwa narodowego (Dz.U. 2022 poz. 835; t.j. Dz.U. 2025 poz. 514 ze zm.).
\BibItem{USyg}{usyg} Ustawa z dnia 14 czerwca 2024 r. o ochronie sygnalistów (Dz.U. 2024 poz. 928).
\BibItem{UODO}{uodo} Ustawa z dnia 10 maja 2018 r. o ochronie danych osobowych (Dz.U. 2018 poz. 1000 ze zm.).
\BibItem{PrLot}{prlot} Ustawa z dnia 3 lipca 2002 r. --- Prawo lotnicze (Dz.U. 2002 nr 130 poz. 1112; t.j. Dz.U. 2025 poz. 1431 ze zm.).
\BibItem{PKD 2025}{pkd} Rozporządzenie Rady Ministrów z dnia 18 grudnia 2024 r. w sprawie Polskiej Klasyfikacji Działalności (PKD) (Dz.U. 2024 poz. 1936).
\BibItem{WT}{wt} Rozporządzenie Rady Ministrów z dnia 17 września 2019 r. w sprawie klasyfikacji rodzajów materiałów wybuchowych, broni, amunicji oraz wyrobów i technologii o przeznaczeniu wojskowym lub policyjnym, na których wytwarzanie lub obrót jest wymagane uzyskanie koncesji (Dz.U. 2019 poz. 1888) --- obejmuje m.in. bezzałogowe statki powietrzne i ich systemy (WT~V pkt~3).
\end{BibList}

\section*{B. Akty prawne Unii Europejskiej}
\begin{BibList}
\BibItem{RODO}{rodo} Rozporządzenie Parlamentu Europejskiego i Rady (UE) 2016/679 z dnia 27 kwietnia 2016 r. w sprawie ochrony osób fizycznych w związku z przetwarzaniem danych osobowych i w sprawie swobodnego przepływu takich danych oraz uchylenia dyrektywy 95/46/WE (ogólne rozporządzenie o ochronie danych) (Dz.Urz. UE L 119 z 4.5.2016, s.~1).
\BibItem{Dual-use}{dualuse} Rozporządzenie Parlamentu Europejskiego i Rady (UE) 2021/821 z dnia 20 maja 2021 r. ustanawiające unijny system kontroli wywozu, pośrednictwa, pomocy technicznej, tranzytu i transferu produktów podwójnego zastosowania (wersja przekształcona) (Dz.Urz. UE L 206 z 11.6.2021, s.~1, ze zm.).
\BibItem{EDF}{edf} Rozporządzenie Parlamentu Europejskiego i Rady (UE) 2021/697 z dnia 29 kwietnia 2021 r. ustanawiające Europejski Fundusz Obronny i uchylające rozporządzenie (UE) 2018/1092 (Dz.Urz. UE L 170 z 12.5.2021, s.~149).
\BibItem{UAS 945}{uas945} Rozporządzenie delegowane Komisji (UE) 2019/945 z dnia 12 marca 2019 r. w sprawie bezzałogowych systemów powietrznych oraz operatorów bezzałogowych systemów powietrznych z państw trzecich (Dz.Urz. UE L 152 z 11.6.2019, s.~1).
\BibItem{UAS 947}{uas947} Rozporządzenie wykonawcze Komisji (UE) 2019/947 z dnia 24 maja 2019 r. w sprawie przepisów i procedur dotyczących eksploatacji bezzałogowych statków powietrznych (Dz.Urz. UE L 152 z 11.6.2019, s.~45).
\BibItem{FDI}{fdi} Rozporządzenie Parlamentu Europejskiego i Rady (UE) 2019/452 z dnia 19 marca 2019 r. ustanawiające ramy monitorowania bezpośrednich inwestycji zagranicznych w Unii --- uchylone z dniem 17 stycznia 2028 r. przez \Zr{fdi2026}.
\BibItem{FDI 2026}{fdi2026} Rozporządzenie Parlamentu Europejskiego i Rady (UE) 2026/1386 z dnia 17 czerwca 2026 r. w sprawie monitorowania inwestycji zagranicznych w Unii i uchylające rozporządzenie (UE) 2019/452 (Dz.Urz. UE L 2026/1386 z 26.6.2026); stosuje się od 17 stycznia 2028 r. (niektóre przepisy od 16 lipca 2026 r.); obejmuje m.in. spółki rozwijające produkty podwójnego zastosowania, wojskowe i drony.
\BibItem{TTBER}{ttber} Rozporządzenie Komisji (UE) 2026/877 z dnia 16 kwietnia 2026 r. w sprawie stosowania art.~101 ust.~3 TFUE do kategorii porozumień o transferze technologii (Dz.Urz. UE L 2026/877 z 21.4.2026); obowiązuje od 1 maja 2026 r. do 30 kwietnia 2038 r.; zastąpiło rozporządzenie (UE) nr 316/2014 --- licencje na technologię, know-how i oprogramowanie.
\BibItem{AI Act}{aiact} Rozporządzenie Parlamentu Europejskiego i Rady (UE) 2024/1689 z dnia 13 czerwca 2024 r. w sprawie ustanowienia zharmonizowanych przepisów dotyczących sztucznej inteligencji (akt w sprawie sztucznej inteligencji) (Dz.Urz. UE L 2024/1689 z 12.7.2024); systemy wyłącznie wojskowe są wyłączone (art.~2 ust.~3), systemy podwójnego zastosowania --- objęte.
\BibItem{Sankcje}{sankcje} Rozporządzenie Rady (UE) nr 833/2014 z dnia 31 lipca 2014 r. dotyczące środków ograniczających w związku z działaniami Rosji destabilizującymi sytuację na Ukrainie (Dz.Urz. UE L 229 z 31.7.2014, s.~1) oraz rozporządzenie Rady (UE) nr 269/2014 z dnia 17 marca 2014 r. (Dz.Urz. UE L 78 z 17.3.2014, s.~6) (środki ograniczające).
\BibItem{TSD}{tsd} Dyrektywa Parlamentu Europejskiego i Rady (UE) 2016/943 z dnia 8 czerwca 2016 r. w sprawie ochrony niejawnego know-how i niejawnych informacji handlowych (tajemnic przedsiębiorstwa) przed ich bezprawnym pozyskiwaniem, wykorzystywaniem i ujawnianiem (Dz.Urz. UE L 157 z 15.6.2016, s.~1).
\BibItem{WBD}{wbd} Dyrektywa Parlamentu Europejskiego i Rady (UE) 2019/1937 z dnia 23 października 2019 r. w sprawie ochrony osób zgłaszających naruszenia prawa Unii.
\end{BibList}

\section*{C. Orzecznictwo}
\begin{BibList}
\BibItem{SN 61/03}{sn6103} Uchwała składu siedmiu sędziów SN z dnia 6 listopada 2003 r., III CZP 61/03, OSNC 2004, nr 5, poz. 69 (zasada prawna) --- kara umowna należy się także wtedy, gdy wierzyciel nie poniósł szkody; dłużnik może żądać jej miarkowania (art.~484 §~2 KC).
\BibItem{SN 304/00}{sn30400} Wyrok SN z dnia 3 października 2000 r., I CKN 304/00, OSNC 2001, nr 4, poz. 59 --- wykorzystanie przez pracownika informacji, co do których przedsiębiorca nie podjął niezbędnych działań w celu zachowania ich poufności, należy traktować jako wykorzystanie powszechnej wiedzy, do której przedsiębiorca nie ma ustawowych uprawnień.
\BibItem{SN 120/07}{sn12007} Wyrok SN z dnia 8 stycznia 2008 r., II PK 120/07 --- odpowiedzialność pracownika za szkodę regulują wyczerpująco przepisy Kodeksu pracy; kara umowna (art.~483 KC) zastrzeżona w stosunku pracy nie może być dochodzona; dopuszczalna jest natomiast w umowie o zakazie konkurencji po ustaniu stosunku pracy (SN powołuje wyrok z 10 października 2003 r., I PK 528/02, OSNP 2004, nr 19, poz. 336).
\BibItem{SN 42/18}{sn4218} Uchwała SN z dnia 7 września 2018 r., III CZP 42/18 (Lex nr 2541937) --- członek zarządu nie może jako pełnomocnik reprezentować spółki z o.o. w umowie z innym członkiem zarządu (art.~210 §~1 KSH); podobnie uchwała SN z 30 stycznia 2019 r., III CZP 71/18. Skutkiem naruszenia jest bezwzględna nieważność czynności (art.~58 §~1 KC) --- tak wyrok SA w Warszawie z 9 kwietnia 2019 r., V ACa 485/18, który obie uchwały przywołuje. Przydatna przy wykładni art.~300$^{79}$ KSH.
\end{BibList}

\section*{D. Literatura}
\begin{BibList}
\BibItem{Malinowski}{malinowski} P. Malinowski (red.), \emph{Prosta spółka akcyjna. Komentarz}, Warszawa 2021.
\BibItem{Opalski}{opalski} A. Opalski, \emph{Prosta spółka akcyjna --- nowy typ spółki handlowej}, cz.~I: „Przegląd Prawa Handlowego” 2019, nr 11, s.~5--15; cz.~II: tamże, nr 12, s.~4--16.
\BibItem{ECFR}{ecfr} \emph{The Simple Joint Stock Company: Emergence of a New Close Company in Poland}, „European Company and Financial Law Review” 2024, t.~21, DOI 10.1515/ecfr-2024-0021.
\BibItem{Jara}{jara} Z. Jara (red.), \emph{Kodeks spółek handlowych. Komentarz}, C.H. Beck (wydanie aktualne).
\BibItem{Gniewek}{gniewek} E. Gniewek, P. Machnikowski (red.), \emph{Kodeks cywilny. Komentarz}, C.H. Beck (wydanie aktualne).
\BibItem{Barta}{barta} J. Barta, R. Markiewicz (red.), \emph{Ustawa o prawie autorskim i prawach pokrewnych. Komentarz}, Wolters Kluwer (wydanie aktualne).
\BibItem{Kostański}{kostanski} P. Kostański (red.), \emph{Prawo własności przemysłowej. Komentarz}, C.H. Beck (wydanie aktualne).
\BibItem{Szwaja}{szwaja} J. Szwaja (red.), \emph{Ustawa o zwalczaniu nieuczciwej konkurencji. Komentarz}, C.H. Beck (wydanie aktualne).
\end{BibList}

\section*{E. Materiały urzędowe}
\begin{BibList}
\BibItem{GUS PKD}{guspkd} Główny Urząd Statystyczny, \emph{Polska Klasyfikacja Działalności PKD 2025} (klasyfikacje.stat.gov.pl).
\BibItem{GUS CPI}{guscpi} Komunikaty Prezesa GUS w sprawie średniorocznego wskaźnika cen towarów i usług konsumpcyjnych ogółem, ogłaszane w Dzienniku Urzędowym RP „Monitor Polski” (podstawa waloryzacji kar i Cennika).
\end{BibList}

\section*{F. Dokumenty Spółki}
\begin{BibList}
\BibItem{US}{us} Umowa spółki Basilisk Systems P.S.A., wersja 3.3-S (\texttt{psa.pdf}; wersja angielska \texttt{psa\_en.pdf}).
\BibItem{UA}{ua} Umowa akcjonariuszy, wersja 1.7 (\texttt{basilisk\_sha.pdf}).
\BibItem{IP}{ip} Umowy przeniesienia praw autorskich (Umowa nr~1 i nr~2), wersja 1.6 (\texttt{basilisk\_ip.pdf}).
\BibItem{Cennik}{cennik} Cennik wynagrodzeń za rozwiązania projektowe (\texttt{basilisk\_cennik.pdf}).
\BibItem{NDA}{nda} Wzór umowy o zachowaniu poufności (\texttt{basilisk\_nda.pdf}).
\BibItem{Polityka}{polityka} Polityka bezpieczeństwa informacji (\texttt{basilisk\_polityka.pdf}).
\BibItem{Podręcznik}{podrecznik} Podręcznik podejmowania decyzji (\texttt{basilisk\_decyzje.pdf}).
\end{BibList}


% ================================================================
\chapter*{Załącznik informacyjny B --- Podstawy prawne per paragraf i ustęp}
\addcontentsline{toc}{chapter}{Załącznik informacyjny B --- Podstawy prawne per paragraf i ustęp}

{\small
\begin{longtable}{@{}>{\raggedright\arraybackslash}p{1.1cm} >{\raggedright\arraybackslash}p{6.0cm} >{\raggedright\arraybackslash}p{7.6cm}@{}}
\toprule
\textbf{Ust.} & \textbf{Podstawa prawna} & \textbf{Orzecznictwo, literatura, komentarz} \\
\midrule
\endfirsthead
\toprule
\textbf{Ust.} & \textbf{Podstawa prawna} & \textbf{Orzecznictwo, literatura, komentarz} \\
\midrule
\endhead
\bottomrule
\endfoot

\Par{firma}{Firma, siedziba i czas trwania}
\Ust{1}{art.~300\U{5} §~1 pkt~1 KSH \Zr{ksh}; art.~43\U{5} §~1--2 KC \Zr{kc} (firma osoby prawnej); znaki towarowe --- art.~120 i nast. PWP \Zr{pwp}}{Firma musi zawierać oznaczenie formy prawnej („prosta spółka akcyjna” albo „P.S.A.”). Ochronę oznaczenia „Basilisk Systems” daje dopiero rejestracja znaku towarowego.}
\Ust{2}{art.~300\U{5} §~1 pkt~1 KSH; art.~41 KC (siedziba = miejscowość)}{W umowie wskazuje się miejscowość, a nie adres; adres zgłasza się do rejestru \Zr{ukrs}.}
\Ust{3}{swoboda umów (art.~353\U{1} KC); \Zr{ukrs} --- oddziały}{Postanowienie fakultatywne.}
\Ust{4}{art.~300\U{5} KSH (postanowienie fakultatywne)}{Brak terminu oznacza czas nieoznaczony.}

\Par{pkd}{Przedmiot działalności i PKD 2025}
\Ust{1}{art.~300\U{5} §~1 pkt~2 KSH}{Opis słowny uzupełnia kody PKD i chroni przed zarzutem działania poza przedmiotem działalności.}
\Ust{2}{\Zr{pkd}; \Zr{guspkd}; \Zr{ukrs}}{Kody spoza ujawnionych przy pierwszym wpisie można zgłosić do rejestru bez zmiany umowy.}
\Ust{3}{koncesja: \Zr{uwt} (art.~7, 10, 17), \Zr{wt}; AI: \Zr{aiact}; kontrola obrotu: \Zr{uots}, \Zr{dualuse}; drony: \Zr{uas945}, \Zr{uas947}, \Zr{prlot}; fundusze obronne: \Zr{edf}}{Wpis kodu PKD nie zastępuje koncesji na wyroby o przeznaczeniu wojskowym ani zezwoleń eksportowych.}

\Par{akcje}{Akcje i kapitał akcyjny}
\Ust{1}{art.~300\U{2} §~3 KSH}{Akcje PSA nie mają wartości nominalnej i nie stanowią części kapitału akcyjnego \Zr{opalski}.}
\Ust{2}{art.~300\U{5} §~1 pkt~3 KSH (liczba, serie i numery akcji)}{Cena emisyjna 0,01 zł jest dopuszczalna; kapitał minimalny wynosi 1 zł (art.~300\U{3} KSH).}
\Ust{3}{art.~300\U{5} §~1 pkt~3 KSH (uprzywilejowanie)}{Uprzywilejowanie musi wynikać z umowy; jego brak ułatwia przyszłą rundę.}
\Ust{4}{art.~300\U{3} §~2 KSH}{Wysokości kapitału akcyjnego nie określa się w umowie; ujawnia ją Rada w rejestrze.}
\Ust{5}{art.~300\U{31} i nast. KSH (rejestr akcjonariuszy)}{Rejestr prowadzi podmiot wybrany przez akcjonariuszy (art.~300\U{31} KSH); ograniczenia i obowiązki związane z akcjami muszą być w nim ujawnione.}
\Ust{6}{art.~300\U{31} i nast. KSH; art.~61 §~2 KC (doręczenie elektroniczne)}{Aktualny adres e-mail warunkuje skuteczną komunikację w trybie §~\ref{art:wz}.}
\Ust{7--8}{art.~300\U{47} KSH (w tym §~3); art.~300\U{15} §~4--6 KSH}{Spółka nie jest nabywcą akcji niezwolnionych --- chroni to przed zarzutem finansowania nabycia własnych akcji.}
\Ust{9}{przepisy KSH o zastawie i użytkowaniu akcji; por. art.~340 KSH (spółka akcyjna)}{Głos zastawnika tylko za zgodą WZ 75\% i po wzmiance w rejestrze.}

\Par{cap-table}{Objęcie akcji i wkłady Założycieli}
\Ust{1--2}{art.~300\U{5} §~1 pkt~3 KSH}{Podział 50/20/10/10/10 dotyczy chwili zawiązania.}
\Ust{3}{art.~300\U{2} §~1--2 KSH; art.~41 i 53 PrAut \Zr{praut}; art.~67 PWP \Zr{pwp}; know-how --- art.~11 UZNK \Zr{uznk}}{Umowa celowo wyłącza wkład w postaci pracy i usług w serii~A oraz sprzęt (sprzedaż albo użyczenie).}
\Ust{4}{art.~5a pkt~34--35, art.~17 ust.~1 pkt~9 i ust.~1e PIT \Zr{pit}}{Część B1 pozwala skorzystać z odroczenia opodatkowania aportu IP przez twórcę.}
\Ust{5}{art.~14 §~1 KSH; art.~300\U{3} §~2 KSH; art.~300\U{9} KSH}{Wkłady pieniężne przed zgłoszeniem, niepieniężne w 14 dni od wpisu.}
\Ust{6}{art.~41 ust.~2 i art.~53 PrAut (forma pisemna, pola eksploatacji); art.~67 PrAut (licencja); art.~33 pkt~9--10 KRO \Zr{kro}; art.~300\U{79} §~1 KSH}{Autorskie prawa majątkowe należą do majątku osobistego twórcy; umowę z dyrektorem podpisuje pełnomocnik \Zr{sn4218}, \Zr{barta}.}
\Ust{7}{---}{Postanowienie informacyjne.}

\Par{emisje}{Emisje akcji i maksymalne liczby akcji}
\Ust{1--3}{art.~300\U{103}--300\U{104} KSH}{Limity ilościowe i termin emisji (do 31.12.2036) chronią akcjonariuszy przed nieograniczonym rozwodnieniem.}
\Ust{4}{przepisy o zmianie umowy PSA (uchwała akcjonariuszy, protokół notarialny, wpis do rejestru) \Zr{ksh}}{Nowa seria inwestorska wymaga zmiany umowy.}
\Ust{5}{art.~300\U{104} KSH}{Rada wykonuje uchwałę emisyjną, nie emituje akcji samodzielnie.}
\Ust{6}{art.~300\U{106} KSH}{Większość 4/5 wynika z ustawy.}
\Ust{7}{art.~300\U{5} §~1 pkt~3 i art.~300\U{104} §~1 pkt~3 KSH; swoboda umów}{Minimalna cena emisyjna.}

\Par{seria-f}{Seria F --- Założyciele Funkcjonalni}
\Ust{1--3}{art.~353\U{1} KC; art.~300\U{104} KSH}{Uchwała kwalifikacyjna nie daje roszczenia o objęcie akcji \Zr{malinowski}.}
\Ust{4}{art.~300\U{104} §~1 pkt~7 KSH; art.~300\U{9} §~1 w zw. z art.~300\U{104} §~2 KSH}{Wkład pracą lub usługami wnosi się w ciągu 3 lat (art.~300\U{9} §~1 stosowany odpowiednio do emisji).}
\Ust{5--7}{art.~300\U{2} KSH; art.~300\U{39} KSH}{Akcje serii~F są zwykłe; podlegają ograniczeniom zbywania i vestingowi.}

\Par{seria-p}{Seria P --- program motywacyjny}
\Ust{1--3}{art.~300\U{103}--300\U{107}, 300\U{114}--300\U{118} (w tym 300\U{116} §~2) i 300\U{106} §~2 KSH}{Warunkowa emisja wymaga wpisu zmiany umowy przed przyznaniem praw.}
\Ust{4--5}{art.~24 ust.~11 i nast. PIT \Zr{pit} (programy motywacyjne)}{Skutki podatkowe dla uczestników wymagają odrębnej analizy przy tworzeniu regulaminu.}

\Par{zbywanie}{Zbywanie i obciążanie akcji}
\Ust{1}{art.~300\U{39} §~1 KSH; art.~300\U{36} KSH (forma dokumentowa); art.~77\U{2} KC; art.~64 KC i art.~1047 KPC \Zr{kpc}}{Milczenie Rady nie jest zgodą; ścieżka uchwały akcjonariuszy i powództwa z art.~64 KC zapobiega blokadzie.}
\Ust{2}{lit.~a: \Zr{uki}, \Zr{fdi}, \Zr{fdi2026}, \Zr{sankcje}, \Zr{usank}, \Zr{uots}; lit.~b: \Zr{uki}; lit.~e: swoboda umów (art.~300\U{39} KSH)}{Zamknięty katalog przyczyn odmowy chroni akcjonariusza przed arbitralnością.}
\Ust{3}{art.~300\U{39} §~2--6 KSH}{Wyłączenie obowiązku wskazania nabywcy, z wyjątkiem odmowy wobec konkurenta.}
\Ust{4}{art.~300\U{39} §~1 KSH; art.~57 §~2 KC}{Lock-up 12 miesięcy --- typowe zabezpieczenie stabilności zespołu.}
\Ust{5}{art.~300\U{40} KSH}{Zbycie akcji niepokrytej wymaga zgody Spółki z mocy ustawy.}
\Ust{6}{art.~353\U{1} KC}{Prawa pierwszeństwa i tag/drag-along wymagają zmiany umowy albo umowy akcjonariuszy \Zr{ua}.}

\Par{obowiazek-zbycia}{Obowiązek zbycia akcji po ustaniu zaangażowania}
\Ust{1--3}{obowiązki związane z akcją przewidziane w umowie (por. art.~356 KSH w spółce akcyjnej); art.~353\U{1} KC}{Vesting w formie obowiązku zbycia --- akcje dają pełne prawa do chwili zbycia \Zr{malinowski}.}
\Ust{4}{art.~393 KC (umowa na rzecz osoby trzeciej); art.~58 i 5 KC}{Cena emisyjna jako cena odkupu wynika z mechanizmu przyjętego przy zawiązaniu; ocena pod kątem zasad współżycia.}
\Ust{5}{art.~353\U{1} KC}{Automatyczne zwolnienie przy dobrym odejściu ogranicza ryzyko sporu.}
\Ust{6}{art.~101 §~1 KC (pełnomocnictwo nieodwołalne); art.~300\U{79} KSH}{Umowa wykonawcza i pełnomocnictwo --- \Zr{ua} §~5.}
\Ust{7}{art.~300\U{31} i nast. KSH}{Obowiązek ujawnia się w rejestrze, aby wiązał nabywców akcji.}
\Ust{8}{art.~922 KC; art.~1035 KC}{Spadkobiercy wstępują w prawa i obowiązki związane z akcjami.}

\Par{rada}{Organy Spółki i Rada Dyrektorów}
\Ust{1--2}{art.~300\U{73} §~1 i §~3 KSH}{System monistyczny \Zr{opalski}, \Zr{ecfr}.}
\Ust{3}{art.~300\U{56} §~2--4 KSH}{Kadencja wspólna w latach obrotowych; limit 6 miesięcy po upływie kadencji to postanowienie umowne, dopuszczalne dzięki zastrzeżeniu „chyba że umowa spółki stanowi inaczej”.}
\Ust{4}{art.~300\U{76} KSH; art.~300\U{75} §~2--3 KSH}{Co najmniej jeden dyrektor niewykonawczy zapewnia funkcję nadzorczą.}
\Ust{5--6}{art.~300\U{75} i 300\U{55} KSH (konflikt interesów); art.~78\U{1} i 77\U{2} KC (forma głosowania)}{Zasady głosowania i konflikt interesów --- dyrektor nie głosuje w sprawie, w której ma sprzeczny interes.}
\Ust{7}{art.~300\U{75} §~2--3 KSH; art.~300\U{81} KSH}{Limity wewnętrzne; nie wpływają na ważność czynności wobec osób trzecich (§~\ref{art:reprezentacja} ust.~4). Progi kwotowe można złagodzić tylko zmianą umowy; uchwała akcjonariuszy może je wewnętrznie zaostrzyć (art.~300\U{53}, 300\U{77} §~2 i 17 §~3 KSH; \Zr{malinowski}).}
\Ust{8}{art.~353\U{1} KC}{Regulamin Rady przyjmuje Rada.}

\Par{reprezentacja}{Reprezentacja Spółki}
\Ust{1}{art.~300\U{78} §~1 KSH}{Reprezentacja łączna (dwóch dyrektorów albo dyrektor z prokurentem).}
\Ust{2}{art.~109\U{1}--109\U{8} KC (prokura); art.~300\U{75} §~3 KSH}{Prokurę może odwołać każdy dyrektor.}
\Ust{3}{art.~300\U{79} §~1, 3 i 4 KSH; art.~15 KSH}{Na tle art.~210 KSH umowa zawarta z naruszeniem szczególnej reprezentacji jest nieważna, a pełnomocnikiem nie może być inny członek organu \Zr{sa48518}, \Zr{jara}; szczególny tryb obejmuje wszystkie czynności związane z umową --- także jej zmianę i rozwiązanie \Zr{sn4218}; rekomendacja: pełnomocnikiem z uchwały akcjonariuszy nie powinien być dyrektor. Droga przez dyrektora niewykonawczego wynika wprost z art.~300\U{79} §~3 KSH.}
\Ust{4}{art.~300\U{77} §~2, art.~300\U{53} i art.~17 §~3 KSH}{Ograniczenia wewnętrzne nie działają wobec osób trzecich.}

\Par{wz}{Uchwały akcjonariuszy --- tryb, zdalne głosowanie i terminy}
\Ust{1}{art.~300\U{80} §~1--2 KSH}{Głosowanie e-mailem nie wymaga zgody wszystkich akcjonariuszy; na piśmie --- wymaga.}
\Ust{2}{art.~77\U{2} KC; art.~61 §~2 KC}{E-mail z adresu z rejestru = forma dokumentowa.}
\Ust{3}{art.~300\U{80} §~1--3 KSH; terminy --- art.~353\U{1} KC}{Tryb poza WZ; terminy 3 dni robocze -- 21 dni wynikają z umowy, nie z ustawy.}
\Ust{4}{art.~300\U{90} KSH}{WZ bez formalnego zwołania.}
\Ust{5--6}{art.~300\U{92} KSH}{Udział zdalny według regulaminu WZ.}
\Ust{7}{art.~300\U{5} KSH; art.~353\U{1} KC}{Większości liczone od wszystkich głosów w Spółce.}
\Ust{8}{art.~300\U{100} KSH}{Księga protokołów.}

\Par{wz-zastrzezone}{Sprawy zastrzeżone dla Walnego Zgromadzenia}
\Ust{1}{art.~300\U{81} KSH; art.~15 §~1 KSH (lit.~j); art.~17 §~1 KSH}{Czynność bez uchwały wymaganej ustawą jest nieważna; bez zgody wymaganej tylko umową --- ważna, z odpowiedzialnością dyrektorów.}
\Ust{2}{art.~300\U{106} KSH}{4/5 przy wyłączeniu prawa poboru.}
\Ust{3}{art.~300\U{98} §~1 KSH}{Sprawy zwykłe --- bezwzględna większość oddanych.}
\Ust{4}{art.~300\U{75} §~2 KSH; art.~300\U{81} pkt~3 KSH}{Uchwała WZ jako zgoda, nie zastępuje uchwały Rady.}

\Par{ip}{Własność intelektualna, poufność i bezpieczeństwo}
\Ust{1}{art.~12, 41 ust.~2--3, 53 i 74 ust.~3 PrAut; art.~11 ust.~3 i art.~22 PWP}{Utwory pracownicze i zamówione; wymóg formy pisemnej; potwierdzenia kwartalne \Zr{barta}, \Zr{kostanski}, \Zr{ip}.}
\Ust{2}{art.~67 PrAut (licencja); art.~65 KC; \Zr{ttber}}{Licencja na technologię sprzed zawiązania na 5 lat z przedłużeniem; przy licencjach między przedsiębiorcami uwzględnić wyłączenie grupowe TTBER (m.in. zakaz wyłącznego zwrotnego przeniesienia ulepszeń licencjobiorcy).}
\Ust{3}{art.~11 UZNK \Zr{uznk}; \Zr{tsd}; art.~100 §~2 pkt~4 KP \Zr{kp}}{Ochrona tajemnicy wymaga realnych działań Spółki \Zr{sn30400}, \Zr{szwaja}, \Zr{polityka}.}
\Ust{4}{\Zr{uots}, \Zr{dualuse}, \Zr{sankcje}, \Zr{usank}}{Transfer technologii także wewnątrz Spółki może wymagać zezwolenia.}
\Ust{5}{art.~353\U{1} KC; dla pracowników art.~101\U{1}--101\U{4} KP}{Zakaz konkurencji w okresie zaangażowania lub funkcji; po odejściu --- \Zr{ua} §~8.}
\Ust{6}{art.~415, 471 KC; art.~730 i nast. KPC}{Kary umowne są w umowie akcjonariuszy \Zr{sn6103}; wobec pracowników niedopuszczalne \Zr{sn12007}.}

\Par{zysk}{Zysk i wypłaty}
\Ust{1}{art.~300\U{5} KSH; art.~353\U{1} KC}{Rekomendacja reinwestowania nie jest zakazem wypłat.}
\Ust{2}{art.~300\U{15} KSH (test wypłacalności)}{Wypłata naruszająca test rodzi odpowiedzialność dyrektorów.}
\Ust{3}{\Zr{uor}; art.~300\U{47} §~2 pkt~3 i §~7 KSH}{Kapitał rezerwowy na akcje własne.}
\Ust{4}{art.~3 ust.~1 pkt~9 UoR \Zr{uor}}{Wydłużony pierwszy rok obrotowy.}

\Par{likwidacja}{Rozwiązanie i likwidacja Spółki}
\Ust{1--2}{przepisy KSH o rozwiązaniu i likwidacji PSA}{Uchwała 75\% wszystkich głosów.}
\Ust{3}{\Zr{uots}, \Zr{dualuse}, \Zr{uoni}}{Klasyfikacja eksportowa przed zbyciem nośników i prototypów.}
\Ust{4}{przepisy KSH o podziale majątku}{Podział proporcjonalny do liczby akcji.}

\Par{final}{Postanowienia końcowe}
\Ust{1}{art.~300\U{98} §~2--3 i art.~300\U{100} §~2 KSH}{Zaostrzenie obowiązków związanych z akcją wymaga zgody akcjonariusza, którego dotyczy.}
\Ust{2}{art.~58 §~3 KC}{Klauzula salwatoryjna.}
\Ust{3}{art.~65 KC}{Wykładnia dynamiczna odesłań.}
\Ust{4}{art.~183\U{1} i nast. KPC (mediacja); art.~730 KPC}{Mediacja nie wstrzymuje terminów zaskarżenia uchwał.}
\Ust{5}{art.~300\U{5} KSH}{Załącznik nr~1 jest częścią umowy; niniejszy komentarz --- nie.}
\Ust{6}{art.~300\U{79} §~1 KSH; art.~300\U{31} KSH}{Pierwsze uchwały w akcie zawiązania.}

\multicolumn{3}{@{}l}{\rule{0pt}{1.4em}\textbf{Załącznik nr~1 do umowy}} \\ \midrule
\Ust{A--C}{art.~300\U{5} §~1 pkt~3 KSH; art.~5a pkt~34--35, art.~17 ust.~1e i art.~22 ust.~1e PIT; art.~74 ust.~3 PrAut}{Część B1 tylko dla twórcy posiadającego prawa; program pracowniczy należy do pracodawcy.}
\end{longtable}
}

\end{document}
